\documentclass{exos}
\usepackage{main}

\date{}
\author{Première Spécialité Mathématiques}

\begin{document}
\title{Définition de la fonction exponentielle : Existence}
\maketitle
\section{Méthode d'Euler et Constante de Neper}
On cherche à déterminer une fonction $f$ qui vérifie
\begin{equation}
\label{Exponentielle}
\begin{cases}
f'(x) = f(x)& \text{ pour tout } x \in \R\\
f(0) = 1
\end{cases}
\end{equation}
Pour déterminer les valeurs que peuvent prendre une telle fonction $f$, nous allons employer la \textbf{méthode d'Euler}. Elle consiste à se souvenir que
\begin{equation*}
f'(x) = \lim_{h \to 0} \dfrac{f(x + h) - f(x)}{h}
\end{equation*}
Alors, en prenant $h$ suffisamment petit, on obtient 
\begin{equation}
\label{Euler}
f'(x) \simeq \dfrac{f(x + h) - f(x)}{h} 
\end{equation}

\subsection{Détermination de $f(1)$}

On se donne une fonction $f$ vérifiant \eqref{Exponentielle}.

\begin{alphaquestions}
\item En utilisant \eqref{Exponentielle} et \eqref{Euler} avec $x = 0$ et $h = 1$, en déduire que

\begin{equation*}
f(0) \simeq \dfrac{f(1) - f(0)}{1}
\end{equation*} 
\item En déduire une première approximation de $f(1)$.
\item On remplace $h = 1$ par $h = 0,5 = \dfrac{1}{2}$. En choisissant $x = 0$ puis $x = 0,5$, utliser \eqref{Exponentielle} et \eqref{Euler} pour montrer les deux approximation suivantes :
\begin{equation*}
\begin{cases}
f(0) &\simeq \dfrac{f(0,5) - f(0)}{0,5}\\
f(0,5) &\simeq \dfrac{f(1) - f(0,5)}{0,5}\\
\end{cases}  
\end{equation*}
\item En déduire une approximation de $f(0,5)$ à l'aide de la première ligne. Puis une approximation de $f(1)$ à l'aide de la deuxième ligne.
\item Montrer que cette approximation de $f(1)$ s'obtient grâce à la formule
\begin{equation*}
f(1) \simeq \left(1 + \dfrac{1}{2}\right)^2
\end{equation*} 
\item Recommencer le même travail avec $h = 0,25 = \dfrac{1}{4}$. Il faudra utiliser la méthode d'Euler quatre fois, avec $x = 0$, $x = 0,25$, $x = 0,5$ et $x = 0,75$. Montrer que la formule permettant d'approcher $f(1)$ dans ce cas est
\begin{equation*}
f(1) \simeq \left(1 + \dfrac{1}{4}\right)^4
\end{equation*}
\item Conjecturer la formule permettant de calculer cette approximation pour $h = 0,1 = \dfrac{1}{10}$.

En passant à la limite $h \to 0$, on obtient que $f(1)$ vaut un nombre appelé \textbf{constante de Neper}, notée $e$. Cette constante vaut environ $2,718\dots$.

\subsection{Détermination de $f(n)$}

La donnée de $f(1)$ va être précieuse pour déterminer le reste des valeurs prises par cette fameuse fonction. Pour cela, on va étudier la fonction $g : x \mapsto \dfrac{f(x + 1)}{f(x)}$. On admet ici que $f$ prend des valeurs strictement positives.

\item Rappeler la formule de dérivation d'un quotient $\dfrac{u}{v}$.
\item En justifiant que $g$ est dérivable sur $\R$, montrer que pour tout $x \in \R$,
\begin{equation*}
g'(x) = \dfrac{f'(x + 1)f(x) - f(x+1)f'(x)}{f(x)^2}
\end{equation*}
\item En déduire que $g$ est constante telle que pour tout $x \in \R$, $g(x) = f(1)$.
\item En déduire que pour tout $x$,
\begin{equation*}
f(x+1) = f(1)f(x)
\end{equation*}
On pose la suite $(u_n)_{n \in \N}$ définie par $u_n = f(n)$ pour $n \in \N$.
\item Justifier que $(u_n)$ est géométrique de raison $e$ et de premier terme $1$.
\item En déduire une expression de $f(n)$ en fonction de $n$.

Ainsi, la fonction $f$ est liée aux puissances de $e$. On peut généraliser cela en notant $f(x)=e^x$, pour n'importe quel réel $x$. La fonction $f$ solution est nommée \textbf{fonction exponentielle}.

\end{alphaquestions}
\newpage
\title{Définition de la fonction exponentielle : Unicité}
\maketitle
\section{Équation différentielle}
On s'intéresse aux fonctions $f : \R \to \R$ qui vérifient \textbf{l'équation différentielle} suivante :
\begin{equation*}
f' = f
\end{equation*}
Autrement dit, pour tout $x \in \R$, si $f$ est solution de cette équation, alors,
\begin{equation*}
f'(x) = f(x)
\end{equation*}
\subsection{Sans condition initiale}
On suppose que la fonction $f$ est une telle solution.
\begin{alphaquestions}
\item Montrer que la fonction $g : x \mapsto 2f(x)$ est définie et dérivable sur $\R$. Montrer que sa dérivée vérifie $g'(x) = g(x)$.
\item Même question avec la fonction $h : x \mapsto 17f(x)$.
\item En déduire qu'il y a une infinité de fonctions $f'$ vérifiant l'équation différentielle $f' = f$.

\subsection{Ajout de la condition initiale}
On suppose alors que $f$ est un telle solution, et qu'en plus, $f$ vérifie :
\begin{equation*}
f(0) = 1
\end{equation*}
\item Soit $h : x \mapsto f(x) \times f(-x)$. Montrer que $h$ est défini et dérivable sur $\R$.
\item Calculer la dérivée de $h$. En déduire que $h$ est une fonction constante.

Cette constante vaut $h(0) = f(0) \times f(- 0) = 1 \times 1 = 1$

\item En déduire que pour tout $x$, $f(x) > 0$.

Ainsi, si $f$ est solution de $f' = f$, et si $f(0) = 1$, $f$ est toujours strictement positive.

\subsection{Unicité}
On suppose maintenant que $f$ et $g$ vérifient tous les deux les même critères :
\begin{itemize}
\item $f' = f$ et $f(0) = 1$
\item $g' = g$ et $g(0) = 1$
\end{itemize}
\item On pose $h(x) = \dfrac{g(x)}{f(x)}$. Montrer que $h$ est définie et dérivable sur $\R$.
\item Calculer la dérivée de $h$, et en déduire que $h$ est constante.
\item Conclure que $g = f$.
\end{alphaquestions}
En conclusion, si une fonction $f$ vérifie $f' = f$ et $f(0) = 1$. Alors cette fonction est unique.

Or, cette fonction existe bel et bien : c'est la \textbf{fonction exponentielle}.

\end{document}